\documentclass{article}

\usepackage{./styles/main_uk}
\graphicspath{{./figure/}}
\addbibresource{src/references.bib}

\begin{document}

% ===================== ABSTRACTS FIRST =====================
% TODO: Replace placeholders with your real topic, goals, results, metrics.
\begin{uaabstract}
  Бакалаврська кваліфікаційна робота присвячена розробленню та експериментальному
  дослідженню гібридного підходу до класифікації зображень рукописного тексту. Мета роботи —
  підвищити точність класифікації у випадку малих і середніх навчальних вибірок за рахунок
  розширення простору ознак виходами ймовірнісної нейронної мережі та подальшого навчання
  класичних алгоритмів машинного навчання.

  Наукова новизна полягає у поєднанні PNN як генератора інформативних ознак із
  лінійними та ансамблевими класифікаторами, що дає змогу зменшити розмірність пошуку
  гіперпараметрів і підвищити стійкість до варіативності почерку. Практична цінність — у створенні
  відтворюваного конвеєра, який дозволяє порівнювати моделі, підбирати гіперпараметри та
  автоматично формувати звіт.

  Експерименти на реальному біомедичному наборі демонструють підвищення F1-міри на
  $5$–$8\,\%$ порівняно з базовою лінією.

  \textbf{Загальний обсяг}: 70 сторінок, 33 рисунки, 34 посилання.

  \textbf{Ключові слова}: ймовірнісна нейронна мережа, класифікація, рукописний текст, розширення
  вибірки, машинне навчання.
\end{uaabstract}

\begin{enabstract}
  This thesis develops and evaluates a hybrid approach to handwritten text image classification.
  The goal is to improve accuracy on small and medium datasets by expanding the feature space with
  outputs of a Probabilistic Neural Network (PNN), followed by training standard machine-learning
  classifiers.

  The contribution is a reproducible pipeline that couples a PNN feature generator with linear
  and ensemble classifiers, reducing hyperparameter search while improving robustness to handwriting
  variation. Experiments on a real biomedical dataset show a $5$–$8\,\%$ F1 improvement over a
  baseline.

  \textbf{Total length}: 70 pages, 33 figures, 34 references.

  \textbf{Keywords}: probabilistic neural network, classification, handwriting, feature expansion, machine learning.
\end{enabstract}

% ===================== TABLE OF CONTENTS =====================
\clearpage
\tableofcontents
\clearpage

% ===================== ABBREVIATIONS =====================
% TODO: Adjust to your topic; keep only what you use in the text.
\section*{ПЕРЕЛІК УМОВНИХ ПОЗНАЧЕНЬ}
\addcontentsline{toc}{section}{ПЕРЕЛІК УМОВНИХ ПОЗНАЧЕНЬ}
\begin{description}
  \item[API] Application Programming Interface --- програмний інтерфейс.
  \item[AI] Artificial Intelligence --- штучний інтелект.
  \item[ML] Machine Learning --- машинне навчання.
  \item[DL] Deep Learning --- глибинне навчання.
  \item[PNN] Probabilistic Neural Network --- ймовірнісна нейронна мережа.
  \item[UMAP] Uniform Manifold Approximation and Projection.
  \item[t-SNE] t-distributed Stochastic Neighbor Embedding.
  \item[CART] Classification and Regression Tree.
\end{description}
\clearpage

% ===================== INTRODUCTION =====================
\section*{ВСТУП}
\addcontentsline{toc}{section}{ВСТУП}
% (Why, What, How, Results, Structure)
Оцифровування документів і рукописів створює попит на точні та відтворювані методи
класифікації текстових зображень. \textbf{Мета роботи} -- підвищити точність класифікації в
умовах обмежених даних.

\textbf{Завдання дослідження:}
\begin{itemize}
  \item виконати огляд і систематизацію методів сегментації й виділення ознак \cite{otsu1979,canny1986};
  \item обґрунтувати гібридний підхід із розширенням ознак через PNN;
  \item розробити конвеєр навчання, валідації та оцінювання якості;
  \item провести експерименти та аналіз результатів на реальних даних;
  \item сформулювати висновки, обмеження та напрями подальших досліджень.
\end{itemize}

Стаття \cite{lecun1998gradient} класично формулює підхід ``градієнтне навчання'' до розпізнавання
документів; методи згорткових мереж \cite{krizhevsky2012imagenet,he2016deep} і оптимізації \cite{kingma2014adam}
становлять технічну основу сучасних рішень. Для візуалізації високорозмірних ознак застосовують
t-SNE \cite{tsne2008} та UMAP \cite{umap2018}.

\textbf{Структура роботи.} Розділ~1 подає огляд і постановку задачі. Розділ~2 описує
математичну модель і алгоритмічні компоненти. Розділ~3 присвячено реалізації та експериментам.
Завершують роботу узагальнюючі висновки.

% ===================== CHAPTER 1 =====================
\section{ЗАГАЛЬНІ ПОЛОЖЕННЯ}

\subsection{Проблемний контекст і постановка задачі}
Ціль -- побудувати класифікатор $f:\mathcal{X}\to\mathcal{Y}$, де $\mathcal{X}$ -- простір зображень
символів/слів, $\mathcal{Y}$ -- множина класів. \textit{Вхідні} дані: зображення в градаціях сірого,
неоднорідні умови сканування, варіативність почерку. \textit{Вихід}: мітка класу та,
за наявності, конфіденси.

\subsection{Огляд предметної області}
Контурні методи й порогування залишаються базовими інструментами попередньої обробки
\cite{otsu1979,canny1986}. Глибинні моделі \cite{lecun1998gradient,krizhevsky2012imagenet,he2016deep}
забезпечують високу якість, але потребують великих вибірок та ретельного налаштування.

\subsubsection{Базові етапи обробки}
\begin{enumerate}
  \item нормалізація та вирівнювання яскравості;
  \item бінаризація (наприклад, Оцу \cite{otsu1979});
  \item виділення контурів (Кенні \cite{canny1986});
  \item сегментація символів/слів; екстракція ознак.
\end{enumerate}

\paragraph{Ілюстративний приклад.}
Схема конвеєра наведена на рис.~\ref{fig:pipeline}.

\begin{figure}[htbp]
  \centering
  \includegraphics[width=0.9\linewidth]{example-figure.png}
  \caption{Схема запропонованого конвеєра класифікації рукописного тексту}
  \label{fig:pipeline}
\end{figure}

\subsection{Формулювання підходу}
Гіпотеза: \emph{виходи PNN як узагальнені апостеріорні оцінки} можуть слугувати компактним і
стабільним представленням для класичних класифікаторів. Отже, шукаємо відображення
$g:\mathcal{X}\to\mathbb{R}^k$ (PNN-ознаки) та класифікатор $h:\mathbb{R}^k\to\mathcal{Y}$, такий що
$f(x)=h(g(x))$.

% ===================== CHAPTER 2 =====================
\section{ІНФОРМАЦІЙНЕ ТА МАТЕМАТИЧНЕ ЗАБЕЗПЕЧЕННЯ}

\subsection{Передобробка та виділення ознак}
Позначимо $I:\Omega\to[0,1]$ — нормоване зображення. Адаптивне порогування
$
  T^\ast=\arg\max_T \ \mathrm{BetweenClassVar}(I,T)
$
реалізує критерій Оцу \cite{otsu1979}.
Градієнтні оператори використовують для детекції границь \cite{canny1986}.

\subsection{Гібридна модель PNN\,$\to$\,класифікатор}
Нехай $\phi_\sigma(x)$ -- вектор виходів PNN із параметром згладжування $\sigma$. Далі
навчаємо $h$ (SVM, Random Forest або логістичну регресію) на $\phi_\sigma(x)$.
Для оптимізації параметрів використовуємо крос-валідацію $k$-fold.
\begin{equation}
  \label{eq:objective}
  \hat{\theta}=\arg\max_{\theta} \ \frac{1}{k}\sum_{i=1}^{k}\mathrm{F1}\bigl(h_{\theta}(\phi_\sigma(X_{-i})),\,Y_{-i}\bigr).
\end{equation}
Рівняння~\eqref{eq:objective} мінімізує \emph{узагальнену помилку} та захищає від
перенавчання.

\paragraph{Метрики якості.}
Використовуються Accuracy, Precision, Recall та F1. Зразок оформлення наведено в табл.~\ref{tab:metrics}.
\begin{table}[htbp]
  \caption{Порівняння метрик якості для трьох підходів}
  \label{tab:metrics}
  \centering
  \begin{tabular}{lccc}
    \hline
    Метод    & Точність & Повнота & F1   \\
    \hline
    Baseline & 0.84     & 0.81    & 0.82 \\
    PNN+SVM  & 0.89     & 0.87    & 0.88 \\
    PNN+RF   & 0.90     & 0.88    & 0.89 \\
    \hline
  \end{tabular}
\end{table}

\subsection{Візуалізація та інтерпретація}
Для інспекції кластерної структури застосовують t-SNE \cite{tsne2008} та UMAP \cite{umap2018}.
\begin{enumerate}
  \item t-SNE добре розкриває локальну структуру, але чутливий до гіперпараметрів;
  \item UMAP частіше зберігає глобальну структуру та масштабується краще.
\end{enumerate}

\subsection{Загальні зауваги щодо відтворюваності}
\begin{itemize}
  \item фіксуйте зерна генераторів випадкових чисел;
  \item зберігайте версії даних і середовища виконання;
  \item логування гіперпараметрів і метрик у єдиному форматі (CSV/JSON).
\end{itemize}

% ===================== CHAPTER 3 =====================
\section{ПРОГРАМНЕ ТА ТЕХНІЧНЕ ЗАБЕЗПЕЧЕННЯ}

\subsection{Архітектура програмної системи}
Система складається з модулів: \textit{io} (завантаження даних), \textit{preprocess},
\textit{features} (PNN), \textit{models} (класифікатори), \textit{eval} (метрики, звіти).
% TODO: у звіті додайте схему класів/модулів вашої реалізації.

\subsection{Середовище виконання та вимоги}
Python~3.x; бібліотеки для обробки зображень і ML; достатня RAM для навчання.
Для оптимізації використовуємо Adam \cite{kingma2014adam} та циклічні графіки навчання за потреби.

\subsection{Експериментальний протокол}
\begin{itemize}
  \item поділ на train/validation/test без перетину екземплярів одного автора;
  \item стратифікована вибірка для збалансованості класів;
  \item $k$-fold крос-валідація для підбору $\sigma$ (PNN) та гіперпараметрів $h$;
  \item статистична перевірка відмінностей (наприклад, парний t-тест на F1).
\end{itemize}

\subsection{Загрози валідності}
Похибки розмітки, дрейф даних, витік тестової інформації через препроцесинг. Необхідно
перевіряти чутливість до параметрів та узагальнюваність на інших корпусах.

% ===================== GENERAL CONCLUSIONS =====================
\section*{ЗАГАЛЬНІ ВИСНОВКИ}
\addcontentsline{toc}{section}{ЗАГАЛЬНІ ВИСНОВКИ}
Запропоновано гібридний підхід PNN\,$\to$\,класифікатор, який підвищує F1 у сценаріях з
обмеженими даними. Розроблено відтворюваний конвеєр експериментів, проведено
порівняльний аналіз і сформульовано практичні рекомендації щодо вибору моделей та
гіперпараметрів.

% ===================== BIBLIOGRAPHY =====================
\printbibliography[heading=bibintoc,title={СПИСОК ВИКОРИСТАНИХ ДЖЕРЕЛ}]

% ===================== APPENDICES =====================
\appendix

\section*{Додаток А}
\addcontentsline{toc}{section}{Додаток А}
% TODO: Перенесіть сюди додаткові таблиці/ілюстрації, приклади даних.

\section*{Додаток Б}
\addcontentsline{toc}{section}{Додаток Б}
% TODO: Посилання на репозиторій, інструкції відтворення (CLI, конфігурація, версії).

\end{document}
